% 付録再編草案・第3部：数列全体と母関数

\section{数列全体を一つの対象にまとめると,何が見えるか}

ここまでは,次の項をどう作るか,隣り合う項の差から何が読めるかを見てきた.
しかし,漸化式が決めているのは一つの項だけではない.
初項から先の全項をつなぐ,数列全体の骨組みである.
その骨組みを一度に持ち運べる形へまとめたら,何が見えるだろう.

数列
\[
  a_0,a_1,a_2,a_3,\ldots
\]
に対して,項の番号を印にして
\[
  A(x)=a_0+a_1x+a_2x^2+a_3x^3+\cdots
  =\sum_{n=0}^{\infty}a_nx^n
\]
と並べる.
これを母関数という.
$x$ を一つ掛けるごとに項の番号が一つ増えるので,\,$x$ は「何番目の項か」を記録する目盛りのような役割を持つ.
最初は,これを収束する関数としてではなく,数列を係数として並べた形式的な記号として扱おう.
母関数の力は,無限個の項を一つに圧縮することよりも,数列の操作を代数の操作に言い換えられることにある.

たとえば $A(x)$ に $x$ を掛けると,
\[
  xA(x)=a_0x+a_1x^2+a_2x^3+\cdots
\]
となる.
数列で言えば,初めに $0$ を置いて,その後の項を一つ右へずらした形である.
また,二つの数列の項を
\[
  c_n=\sum_{k=0}^{n}a_kb_{n-k}
\]
のように組み合わせる操作は畳み込みと呼ばれ,母関数を掛けたときの $x^n$ の係数になる.
一つの対象を作る方法を二つに分ける,あるいは二つの部品を組み合わせるとき,この畳み込みが自然に現れる.
つまり母関数では,「一つ先へずらす」は $x$ 倍に,「二つの選択を組み合わせる」は積に翻訳される.

この翻訳をフィボナッチ数列で追ってみよう.
\[
  F_0=0,\quad F_1=1,\qquad F_{n+2}=F_{n+1}+F_n
\]
とし,\,$F(x)=\sum_{n=0}^{\infty}F_nx^n$ と置く.
左辺の初めの項を取り出すと,
\[
  F(x)-x=\sum_{n=0}^{\infty}F_{n+2}x^{n+2}
\]
である.
漸化式 $F_{n+2}=F_{n+1}+F_n$ を代入すると,右辺の二つの和はそれぞれ $xF(x)$ と $x^2F(x)$ になる.
したがって
\[
  F(x)-x=xF(x)+x^2F(x)
\]
であり,
\[
  F(x)=\frac{x}{1-x-x^2}
\]
を得る.
項を一つずつ計算する代わりに,数列全体が満たす一つの代数方程式を解いた.
漸化式は母関数の分母に姿を変えたのである.

ここで分母を特性方程式と比べてみる.
\[
  1-x-x^2=0
\]
に対して $r=1/x$ と置くと,\,$r^2-r-1=0$ となる.
つまり母関数の分母の零点は,特性方程式の根の逆数である.
これは偶然ではない.
特性方程式は「一歩ごとの倍率」を調べ,母関数は「全項を並べたときの係数関係」を調べている.
行列の固有値,特性根,母関数の分母は,同じ更新規則を違う入口から見たものだった.

フィボナッチ数列をもう少し遠くから眺めると,長期的な増え方も見えてくる.
特性根は
\[
  \varphi=\frac{1+\sqrt5}{2},
  \qquad
  \psi=\frac{1-\sqrt5}{2}
\]
であり,\,$|\psi|<1<\varphi$ である.
一般項
\[
  F_n=\frac{\varphi^n-\psi^n}{\sqrt5}
\]
では,\,$\psi^n$ の影響は次第に小さくなる.
したがって
\[
  F_n\sim\frac{\varphi^n}{\sqrt5}
\]
となる.
ここで $\sim$ は,二つの量の比が $1$ に近づくという意味である.
大きな $n$ でどの根が勝つかを見れば,正確な一般項を毎回計算しなくても成長の景色が分かる.
母関数の分母からも同じ特性根が取り出されるため,固有値・一般項・母関数・漸近的な増大は一つの話につながる.

母関数が本領を発揮するのは,項どうしの組合せが漸化式に現れるときである.
カタラン数は,たとえば括弧の付け方や二分木の形を数える数列で,
\[
  C_0=1,\qquad C_{n+1}=\sum_{k=0}^{n}C_kC_{n-k}
\]
を満たす.
全体を左側の大きさ $k$ と右側の大きさ $n-k$ に分けると,左の形と右の形を一つずつ選ぶので積 $C_kC_{n-k}$ が現れる.
すべての分け方を足すため,右辺は畳み込みになっている.
母関数 $C(x)=\sum_{n=0}^{\infty}C_nx^n$ を作ると,畳み込みは積に変わり,
\[
  C(x)=1+xC(x)^2
\]
となる.
項ごとには複雑に見えた再帰が,母関数についての二次方程式へ変わった.
これを解くと
\[
  C(x)=\frac{1-\sqrt{1-4x}}{2x}
\]
である.
もう一方の符号では $x=0$ の近くで値が発散してしまい,定数項 $C_0=1$ を持つ母関数にならないため,こちらの枝を選ぶ.
母関数を使うと,非線形な漸化式を力ずくで代数方程式へ移し,そこから項の情報を引き出せる.

この発想は,チェビシェフ多項式にも戻ってくる.
その漸化式に母関数を作ると,
\[
  \sum_{n=0}^{\infty}T_n(x)t^n
  =\frac{1-xt}{1-2xt+t^2}
\]
となる.
これは,漸化式に対応する項を一つずつ並べ,同じ次数の係数を整理すれば得られる.
三角関数で定義した多項式列を,今度は一つの有理関数にまとめている.
漸化式,三角関数,母関数は別々の定義ではなく,同じ列を引き出すための三つの窓になっている.
このように,一つの対象に複数の表現を持たせると,見えにくかった性質が別の窓から見えてくる.

さらに遠くには,「数列を満たす関数を作る」話も広がっている.
階乗は整数 $n$ で $(n+1)!=(n+1)n!$ を満たすが,階乗を整数以外でも使いたい場面がある.
その拡張の一つがガンマ関数で,正の $x$ に対して
\[
  \Gamma(x)=\int_0^\infty t^{x-1}e^{-t}\,dt
\]
と定められる.
部分積分をすると
\[
  \Gamma(x+1)=x\Gamma(x)
\]
が現れ,整数では $\Gamma(n+1)=n!$ となる.
数列の外へ関数を作ると,階乗が積分や確率に現れる新しい景色へつながる.
本節ではその理論を追わないが,チェビシェフ多項式と同じく,漸化式が関数を生み出す一つの例として見ておこう.

確率に目を向けると, 状態の移り変わりも漸化式で記録できる.
晴れ・雨の順に確率を並べた列ベクトルを $\boldsymbol p_n$ とし, 晴れの翌日に晴れる確率を $0.8$, 雨の翌日に晴れる確率を $0.3$ とすると,
\[
  \boldsymbol p_{n+1}
  =
  \begin{pmatrix}0.8&0.3\\0.2&0.7\end{pmatrix}
  \boldsymbol p_n
  =P\boldsymbol p_n
\]
と書ける.
各列の和は $1$ なので, どの状態から出発しても次の確率の合計は $1$ に保たれる.
長く更新したときの定常分布は $P\boldsymbol p=\boldsymbol p$ を満たし, どの初期状態からそこへ近づくかは固有値で調べられる.
線形代数で見た状態の更新が, 確率の予測にもそのまま現れている.

数論では, ユークリッドの互除法が余りを次々に作る更新として表せる.
$a_{n-1}$ を $a_n$ で割った商を $q_n$, 余りを $a_{n+1}$ とすれば,
\[
  a_{n-1}=q_na_n+a_{n+1},
  \qquad 0\le a_{n+1}<a_n
\]
となる.
正の余りは毎回小さくなるため, この計算は有限回で止まり, 最大公約数が求まる.
ここでは一般項を求めず, 更新のたびに小さくなるという性質からアルゴリズムの停止を示した.

同じ数論の見方として, 整数係数が一定の二階漸化式を $m$ で割った余りで考えてみよう.
状態 $(a_n,a_{n+1})$ の各成分は $0,1,\ldots,m-1$ のいずれかなので, 状態は高々 $m^2$ 通りしかない.
次の状態が今の状態だけで決まるなら, いつか同じ状態が現れ, その後は同じ並びを繰り返す.
したがって, 余りの列は周期的になる.
たとえばフィボナッチ数列を $2$ で割った余りは
\[
  0,1,1,0,1,1,0,1,1,\ldots
\]
と周期 $3$ で繰り返す.
実数としての増大や収束を調べるときと, 余りとしての周期を調べるときでは, 同じ更新規則から見えてくる性質が変わる.

一般項を簡単な式で書きにくい数列にも, 調べられることは残っている.
たとえば
\[
  a_0=0,\qquad a_{n+1}=a_n^2+1
\]
では $0,1,2,5,26,\ldots$ と増える.
$n\ge2$ では $a_n\ge2$ であり, $a_{n+1}\ge a_n^2$ だから, 帰納的に
\[
  a_n\ge 2^{\,2^{n-2}}
\]
となる.
一般項を求めなくても, 増加が非常に速いことは分かる.
また, 数列が単調で有界と分かれば収束し, 更新が $a_{n+1}=f(a_n)$ で $f$ が連続なら, 極限 $L$ は $L=f(L)$ を満たす.
単調性, 有界性, 不動点, 増加速度は, 一般項とは別の問いに答えてくれる.

本編では, 漸化式から一般項を取り出す方法を整理した.
付録では, 同じ更新規則を行列で表し, 一歩の差を細かく見て微積分へつなぎ, 全項を母関数にまとめた.
確率の定常分布, 整数の余りの周期, 一般項を使わない増加の評価も, 見たい性質に応じて更新規則を読み替えた例である.
漸化式を理解する道具は, 一般項を得るためだけでなく, その数列がどのように動き, 何を数え, どのような性質を保つかを知るために広がっている.
